\clearpage
\VerdierPhysicalPageAnchor{0244}%
\begingroup\small\setlength{\parskip}{0pt}\setlength{\emergencystretch}{1.5em}
\setlength{\abovedisplayskip}{4pt plus 2pt minus 2pt}%
\setlength{\belowdisplayskip}{4pt plus 2pt minus 2pt}%

\VerdierSemanticAnchor{ch4d-4-5-4-proof-e-exact-diagram}
On en déduit un diagramme commutatif, où les lignes sont exactes :
\[
\begin{tikzcd}[column sep=2.7em,row sep=2.8em]
F_1^{p+r-1}(Y) \arrow[r] \arrow[d] &
F_1^p(L^\bullet) \arrow[r] \arrow[d,"{m(L^\bullet[p])}"'] &
F_1^p(X) \arrow[r] \arrow[d,"{m(X[p])}"'] &
F_1^{p+r}(Y) \arrow[d] \\
F_2^{p+r-1}(Y) \arrow[r] &
F_2^p(L^\bullet) \arrow[r] &
F_2^p(X) \arrow[r] &
F_2^{p+r}(Y).
\end{tikzcd}
\]
L'entier $r$ a été choisi de telle façon que les objets $F_i^q(Y)$ soient
nuls, pour $i=1,2$ et $q\geq p+r-1$, et, par suite, tout revient à montrer
que $m(L^\bullet)$ est un isomorphisme, pour tout objet $L^\bullet$ de
$\operatorname{comp}^{b}(\mathcal A)$ dont les composants sont des éléments
de $M$.

\VerdierSemanticAnchor{ch4d-4-5-4-proof-e-spectral-sequence}
Or le morphisme $m$ induit un morphisme entre les premières suites
spectrales des foncteurs $F_1$ et $F_2$ relatives au complexe $L^\bullet$.
Sur les termes initiaux, ce morphisme :
\[
m_1^{p,q}(L^\bullet)\colon F_1^q(L^p)\longrightarrow F_2^q(L^p),
\]
est, par hypothèse, un isomorphisme. Ces suites spectrales étant
stationnaires (4.5.3), le morphisme $m$ induit un isomorphisme sur les
aboutissements (chap.~II, 4.4.5) et, par suite (4.4.2), le morphisme
$m(L^\bullet)$ est un isomorphisme.

\VerdierSemanticAnchor{ch4d-4-5-4-proof-end}
Ceci achève la démonstration de l'assertion e). L'assertion f) se déduit de
l'assertion e), en passant aux catégories opposées. Le théorème est
démontré.

\VerdierSemanticAnchor{ch4d-4-6}
\medskip\noindent\textbf{4.6. Objet spectral de Cartan-Eilenberg.}

\VerdierSemanticAnchor{ch4d-4-6-1}
\noindent\textbf{Définition 4.6.1.}\quad
On dit qu'un complexe $I^\bullet$ de $\mathcal A$ est un \emph{complexe
injectif de Cartan-Eilenberg} si, pour tout entier $i$, les objets
$\operatorname{Im}(d_{I^\bullet}^{i})$ et $H^i(I^\bullet)$ sont injectifs.

\VerdierSemanticAnchor{ch4d-4-6-2}
\noindent\textbf{4.6.2.}\quad
On notera qu'un complexe injectif de Cartan-Eilenberg $I^\bullet$ est de
type injectif (1.1.5), i.e.\ que ses composants sont des objets injectifs de
$\mathcal A$ et que, pour tout entier $i$, l'objet
$\operatorname{Ker}(d_{I^\bullet}^{i})$ est un objet injectif de
$\mathcal A$. Ceci se vérifie immédiatement, en utilisant les suites exactes
à trois termes qui relient ces différents objets.

\VerdierSemanticAnchor{ch4d-4-6-3-a}
\medskip\noindent\textbf{Proposition 4.6.3.}\quad
\textit{Soit $I^\bullet$ un complexe injectif de Cartan-Eilenberg.}

\noindent\textit{a) Pour tout complexe $Y$, le morphisme canonique :}
\[
\operatorname{Hom}_{K(\mathcal A)}(Y,I^\bullet)
\longrightarrow
\prod_{p\in\mathbb Z}
\operatorname{Hom}_{\mathcal A}\bigl(H^p(Y),H^p(I^\bullet)\bigr)
\]
\textit{est un isomorphisme.}
\endgroup
